% !PRINTED 410
\weberpage{410}
%<<CONT
and the sum $\overset{\mathfrak a}{\Sigma}\dfrac{s-1}{(N\mathfrak a)^s}$, extended over a class $A$, has for $s=1$ a definite limiting value different from zero, which by § 111, 5. is the same for all classes. From this it follows by 2.:
\weq{3}{\overset{A}{\Sigma}(\delta,A)=0,}
except when $\delta=1$ or $=\Delta$,
\weq{4}{\overset{A}{\Sigma}(\delta,A)=h,}
when $\delta=1$ or $=\Delta$ [because $(\Delta,A)=1$, § 104, (10)].

Further, for every class $A_0$ of the principal genus $(\delta,A_0)=1$, while for every class $A$ which does not belong to the principal genus there is at least one $\delta_1$ such that $(\delta_1,A)=-1$. From this it follows, if $2^\nu$ is the number of fundamental divisors,
\weq{5}{\begin{aligned}
\overset{\delta}{\Sigma}(\delta,A_0) &= 2^{\nu},\\
(\delta_1,A)\overset{\delta}{\Sigma}(\delta,A) &= \overset{\delta}{\Sigma}(\delta\delta_1,A),
\end{aligned}}
and since $\delta\delta_1$ runs through the totality of the fundamental divisors together with $\delta$ (if one reduces $\delta\delta_1$ according to the symbolic multiplication in § 104):
\[
(\delta_1,A)\overset{\delta}{\Sigma}(\delta,A) = \overset{\delta}{\Sigma}(\delta,A),
\]
and hence, when $(\delta_1,A)=-1$:
\weq{6}{\overset{\delta}{\Sigma}(\delta,A)=0.}

Accordingly we can determine the double sum
\[
\overset{\delta}{\Sigma}\,\overset{A}{\Sigma}(\delta,A)
\]
in two ways. By (3), (4) there results:
\weq{7}{\overset{\delta}{\Sigma}\,\overset{A}{\Sigma}(\delta,A)=2h,}
and by (5), (6):
\weq{8}{\overset{\delta}{\Sigma}\,\overset{A}{\Sigma}(\delta,A)=2^{\nu}g,}
where $g$ denotes the number of classes of the principal genus (and consequently of every genus). The comparison gives:
\weq{9}{h=2^{\nu-1}g,}
and therefore the theorem:

6. \gesperrt{The number of genera is exactly equal to $2^{\nu-1}$, i.\,e. equal to $\tfrac{1}{2}(Z,R)$, § 108, 9.}

\begin{center}\rule{0.32\linewidth}{0.4pt}\end{center}
